\addcontentsline{toc}{chapter}{Further Research}
\chapter*{Further Research}

This research is an initial parsing of the landscape; it leaves much unexamined. Further work should consider frameworks for the implementation and design of DIs themselves. Legal accounts of the necessary prior justifications and frameworks are a priority - they exist, but insufficiently to provide law-makers with locally-appropriate legislation-oriented roadmaps. This is especially true at the global level, and in regions of the Global South. Economic analyses to support appropriate pricing structures for licensing data will involve the unenviable task of considering how to appropriately value data and DIs' governance services. Political-economy investigations into conditions under which DIs succeed/fail, conditional on a variety of adjustable institutional and contextual parameters, would support evidence-based institutional design. Similarly, securing financial independence and protection against corporate or government capture remains under-considered, especially as core proposed functions (such as verification and data collection) are resource intensive. 

At a technical level, pressing open questions revolve around how to assert control and stewarding rights over data currently generated on privately-provisioned platform infrastructure but which might reasonably be considered commons or private data. Similarly, blockchain and cryptographic tools to develop trust and enforcement are promising areas - although purely technical solutions to human problems should be implemented cautiously.

Across all these areas, mapping existing data ecosystems to better understand the contexts into which DIs might be inserted - or already are - remains vital groundwork. 

